\section{Methodology}
\label{sec:method}

\subsection{Corpus}

The corpus holds \NumPapers{} documents, from which \NumSources{} carry a formal
claim and enter the classification. A work is included when three conditions hold.
It concerns learning or something a learning process produced, meaning a model, its
training data, its parameters, its gradients or its outputs. It makes a formal
claim, meaning a theorem, a definition, or a measurement designed to be comparable
across papers. And somebody's resources bear on that claim, in the sense that
deleting every resource restriction would change what the claim says.

The third condition admits two bodies of work with no machine learning content at
all. Fine-grained cryptography enters because it is the only place where the
comparison between an adversary's work and an honest party's is a defined object.
The bounded-storage literature enters on the same grounds, as the formal home of an
adversary limited in space. Both enter as destinations rather than as subject
matter, and Section~\ref{sec:findings-tables} reports that nothing in the machine
learning half of the corpus has travelled to either.

\subsection{Codebook}

Each work is classified on thirteen fields against a fixed vocabulary, so that two
readers assign the same value. Six fields record the attack scenario and threat
model, and seven record what the work assumes about resources.

The vocabulary was revised four times, each revision driven by a classifier
reporting a work it could not place. The revisions are recorded rather than
smoothed away, because their pattern is itself a result. Two of them added values
for bounds that are not computational, namely a perturbation radius, a corruption
fraction and a privacy budget, without which a work that bounds its adversary
non-computationally would have been recorded as bounding nothing. One separated a
second bounded party from one party bounded under two incomparable resource
notions. Eight of the fourteen gaps reported across the whole pass fell in the
single field recording the attack family, which Section~\ref{sec:findings-tables}
returns to.

\subsection{Procedure}

Each work was read first through its catalogue entry, which records the quantifier
of its main claim and the hypothesis that claim rests on, and then through the
primary document wherever that entry was silent on a field. Four rules decide the
cases where a work admits more than one reading. The claim the abstract leads with
decides the row. The constrained party is whoever the bound binds rather than
whoever the paper is about. When two parties are bounded, the computational bound
decides the row and the other party's unit is recorded beside it. A work that
constrains nobody is recorded as constraining nobody, which is an answer and not a
gap.

Every assignment carries the sentence that produced it. For each work the record
holds a verbatim fragment from the source that determines the constrained party,
the unit and the tightness, together with the definition, theorem or section it
came from. Where a work states no bound, the record holds the claim of absence and
the places that were checked before making it, because an unsupported claim of
absence is worth nothing.

\subsection{Verification}

Every quotation was checked against the primary document by extracting its text
and matching a punctuation-independent fragment. Two passes were run. The first,
over the \NumSources{} works then in the corpus, asked whether a quotation
appeared anywhere in the corpus, and found none absent. That question is weaker
than the failure it was meant to catch. A second pass binds each work to its own
document by title and asks whether the quotation appears in \emph{that} document.
It found four works carrying a neighbour's quotation, an off-by-one rotation
within one batch, which the first pass had passed. All four were re-extracted from
their own sources and verified. Of the works the second pass can bind to a
document, 57 verify and none carries another paper's quotation. A further 85 could
not be bound by title, and 25 record an evidenced claim of absence rather than a
quotation.

Eleven entries failed an automated check and all eleven were false alarms, six from
fragments spanning a line break, two from an error in the checker, one from a claim
of absence written as prose, and two from mathematical notation. A single automated
pass is not sufficient evidence that a quotation is absent, and a pass that asks a
weaker question than the failure mode is worse than none, because it returns a clean
report.

\subsection{Threats to validity}

Three limits bear on the counts reported in
Section~\ref{sec:findings-tables}. Each work was classified once rather than by two
independent readers, so the agreement rate between readers is unmeasured, and the
fixed vocabulary and decision rules are what stand in for that measurement.

One field was re-checked in full before being reported. The column recording
whether a work states what an adversary achieves from public knowledge alone was
applied under two readings during the first pass, one asking for the success rate
available without access to the target and one counting any explicit comparison
quantity. Every row was re-read under the narrower definition and
\BaselineChanged{} changed value. Where two readers checked overlapping rows they
agreed on 24 of 27, which is the only inter-reader agreement figure available here.

The field recording the attack family absorbed eight of the fourteen vocabulary
gaps. That field was built around attacks on a deployed model, and the corpus also
contains attacks on a cryptographic primitive, on a communication channel, on a
training distribution and on an agent's control flow. Its categories are therefore
the least reliable of the thirteen, and the grouping used in
Figure~\ref{fig:sankey} collapses them for that reason.
